\section{Exterior algebra and the square of the interaction}
\label{sec:algebra}
Our first task is to isolate the two-, three-, and four-body terms of $H_Q^2$. The lifting notation below keeps this calculation independent of particle number.

Throughout the operator argument take $Q\ge2$; statements involving limits hold for all sufficiently large integer $Q$. We use the convention that wedges of distinct increasing orthonormal basis vectors have norm one. On
\[
 \FF_Q=\bigoplus_{n=0}^{Q+1}\bigwedge^n U_Q,
\]
let $B(v)$ be the adjoint of left wedge multiplication by $v$. Thus $B(v)$ is conjugate-linear in $v$, and
\[
 B(v\wedge w)=B(w)B(v).
\]
For an orthonormal basis $e_0,\ldots,e_Q$ of $U_Q$, write $c_j=B(e_j)$. These operators satisfy the canonical anticommutation relations, and $\norm{c_j}\le1$.

For a matrix on $\bigwedge^k U_Q$, define its lift by linear extension of
\begin{equation}\label{eq:lift}
 \LL_k(\ket v\bra w)=B(v)^\dagger B(w).
\end{equation}
This map preserves positivity, intertwines rotations, acts as the original matrix on $k$ particles, and vanishes on fewer than $k$ particles. Positivity follows by decomposing a positive matrix into a sum of rank-one positive matrices.

Put $a=Q-1$ and let $V_Q\subset\bigwedge^2 U_Q$ be the spin-$a$ summand. If $v_0,\ldots,v_{2Q-2}$ is an orthonormal basis of $V_Q$, set
\begin{equation}\label{eq:fockH}
 B_p=B(v_p),\qquad H_Q=\LL_2(\Pi_{V_Q})=\sum_{p=0}^{2Q-2}B_p^\dagger B_p.
\end{equation}
This pair-annihilator form is standard in the algebraic treatment of
pseudopotentials~\cite[Sections~II.B--D]{OrtizNussinovDukelskySeidel};
see also \cite[Equation~(2)]{ChenSeidel2015}. We verify its coefficient-one
normalization here. Its restriction to $\bigwedge^N U_Q$ is \eqref{eq:physicalH}. Indeed contraction in two specified tensor positions corresponds to $B(v)/\sqrt{\binom N2}$ under the isometric identification of wedge and antisymmetric tensor spaces; summing over the $\binom N2$ unordered pairs gives \eqref{eq:fockH}.

Define intertwining wedge maps
\[
 W_3:V_Q\otimes U_Q\longrightarrow\bigwedge^3U_Q,
 \qquad W_4:V_Q\otimes V_Q\longrightarrow\bigwedge^4U_Q,
\]
by $W_3(v\otimes u)=v\wedge u$ and $W_4(v\otimes w)=v\wedge w$. We suppress the subscript $Q$ when no confusion can result.

A finite normal-ordering calculation in the sense of Wick~\cite{Wick}
now determines the square. Only contractions between the two middle
pair operators are needed.

\begin{lemma}[Normal-ordered square]\label{lem:square}
On $\FF_Q$,
\begin{equation}\label{eq:square}
 H^2=H+\LL_3\!\left(W_3(W_3^\dagger W_3-I)W_3^\dagger\right)
       +\LL_4(W_4W_4^\dagger).
\end{equation}
\end{lemma}
\begin{proof}
In normal ordering $B_p^\dagger(B_pB_r^\dagger)B_r$, only the two inner annihilators cross the two inner creators. Two, one, or zero contractions leave two-, three-, or four-body terms, respectively. On two particles, $H=\Pi_{V_Q}$ is a projection, so the two-body coefficient is $\Pi_{V_Q}$. The term with no contractions is
\[
 \sum_{p,r}B(v_p\wedge v_r)^\dagger B(v_p\wedge v_r)
 =\LL_4(W_4W_4^\dagger).
\]
There is no extra factor of two: this sum, and the defining basis of $V_Q\otimes V_Q$, both use ordered pairs. On three particles, $H=W_3W_3^\dagger$, since $W_3$ restricted to $v_p\otimes U_Q$ is the map $B_p^\dagger$ from one to three particles. Subtracting its two-body contribution from $H^2$ therefore determines the remaining coefficient as
\[
 (W_3W_3^\dagger)^2-W_3W_3^\dagger
 =W_3(W_3^\dagger W_3-I)W_3^\dagger.
\]
Normal-ordered coefficients are determined successively by these restrictions, proving the identity on Fock space.
\end{proof}
